\documentclass[12pt]{article}
\usepackage[utf8]{inputenc}
\usepackage[spanish]{babel}
\usepackage{amsmath}
\usepackage{amsfonts}
\usepackage{amssymb}
\usepackage{geometry}
\geometry{letterpaper, margin=2.5cm}

\begin{document}

\begin{center}
    \Large \textbf{Evaluación de Examen 1: Unidad I} \\
    \large Física para Ciencias de la Salud (005-1713) \\
    \vspace{0.5cm}
    \textbf{Estudiante:} Bárbara Rodriguez \quad \textbf{C.I.:} 32.640.760
    \vspace{0.3cm} \\
    \large \textbf{Calificación Final: 9/10 puntos}
\end{center}

\vspace{0.5cm}

\section*{Análisis de Resultados}

\subsection*{1. Ángulo plano (Conversión a radianes)}
\textbf{Procedimiento y resultado:} Aplica correctamente el factor de conversión $\left(\frac{\pi \text{ rad}}{180^\circ}\right)$. Simplifica metódicamente la fracción hasta $\frac{5\pi}{12}$ rad, y además calcula su valor decimal equivalente ($1,31$ rad). Todo el procedimiento es correcto y claro. \\
\textbf{Veredicto:} \textbf{Correcto} (2/2 pts).

\subsection*{2. Sistemas de coordenadas (Longitud del hueso)}
\textbf{Procedimiento y resultado:} Utiliza la fórmula formal de la distancia entre dos puntos. Sustituye apropiadamente las coordenadas dadas y resuelve las operaciones algebraicas ($\sqrt{36 + 64}$) de manera impecable para llegar a $10$ cm. Acompaña el procedimiento con un gráfico muy bien elaborado y representativo. \\
\textbf{Veredicto:} \textbf{Correcto} (2/2 pts).

\subsection*{3. Vectores (Fuerza resultante)}
\textbf{Procedimiento y resultado:} La estudiante emplea una estrategia muy ordenada: calcula de forma separada la suma de los componentes $\hat{i}$ ($45 - 15 = 30$) y la de los componentes $\hat{j}$ ($25 + 35 = 60$). Esto le permite ensamblar el vector resultante exacto y sin errores algebraicos: $\vec{F}_R = 30\hat{i} + 60\hat{j}$ N. \\
\textbf{Veredicto:} \textbf{Correcto} (2/2 pts).

\subsection*{4. Potencias de diez (Notación científica y prefijo)}
\textbf{Procedimiento y resultado:} Transforma la medida correctamente a notación científica estándar con un solo dígito entero ($1,25 \times 10^{-5}$). Sin embargo, el enunciado también solicitaba expresar el prefijo adecuado del Sistema Internacional. Al no haber expresado el valor en función de una potencia múltiplo de 3 (como $10^{-6}$), se le dificultó identificar el prefijo (micrómetros), omitiendo esta parte de la respuesta. \\
\textbf{Veredicto:} \textbf{Incompleto / Parcialmente correcto} (1/2 pts).

\subsection*{5. Sistemas de unidades (Conversión de densidad)}
\textbf{Procedimiento y resultado:} Plantea los factores de conversión en cadena de manera perfecta: $\left(\frac{1 \text{ kg}}{1000 \text{ g}}\right)$ y $\left(\frac{1.000.000 \text{ cm}^3}{1 \text{ m}^3}\right)$. Simplifica correctamente los factores numéricos y halla de forma exacta el valor de $1030 \text{ kg/m}^3$. \\
\textbf{Veredicto:} \textbf{Correcto} (2/2 pts).

\vspace{0.5cm}
\section*{Comentarios Generales y Sugerencias}
El examen de la estudiante es excelente. Muestra un desarrollo sumamente ordenado, un manejo pulcro del álgebra y un nivel de detalle destacable en sus procedimientos (como el desglose de los vectores y el gráfico en coordenadas).
\begin{itemize}
    \item Como única sugerencia, se recomienda a la estudiante leer atentamente cada pregunta para asegurarse de cumplir con todos los requerimientos del enunciado (como ocurrió al omitir la identificación del prefijo del S.I. en la pregunta 4).
\end{itemize}

\end{document}